\setlength{\parindent}{0em}
\pagenumbering{gobble}

\section*{Resumen}

El trabajo presenta el diseño, implementación y evaluación de un sistema de trading algorítmico basado en técnicas de Machine Learning para la
identificación y optimización de movimientos LONG/SHORT aplicados a futuros financieros. El objetivo principal es operar de forma autónoma una
cartera diversificada (valor objetivo de 1 millón USD) sobre 13 contratos de futuros (divisas, índices, materias primas y bonos), maximizando
beneficio y controlando riesgo. La arquitectura propuesta es híbrida: combina modelos supervisados (LSTM y GRU) para generar señales direccionales y
una capa de regresión para estimar retornos.

El flujo de datos sigue un pipeline ETL organizado en capas Bronce-Plata-Oro: ingestión de datos OHLCV minuto a minuto, limpieza y estandarización
temporal a UTC, resampleo a velas de 4 horas y cálculo de un amplio conjunto de features (momentum, volatilidad, volumen, estructurales y derivados de
swing charts). Se aplica selección de características (correlación, LASSO) para reducir dimensionalidad y evitar data leakage. La etiqueta objetivo se
construye con el método de triple barrera (Stop Loss, Take Profit, vencimiento temporal) y se complementa con la regresión del retorno logarítmico.

Los modelos se entrenan con PyTorch usando una estrategia walk-forward (moving window) para respetar la estructura temporal y evitar fugas. La
solución consta de tres niveles: L1 (detección de movimiento significativo), L2 (dirección BUY/SELL) y L3 (regresión de retorno); la decisión final se
obtiene por soft-voting ponderado por el retorno estimado. Se emplea RobustScaler y técnicas de regularización (dropout) para garantizar la robustez
frente a outliers y cambios de régimen.

El backtest se realiza sobre datos de 4 horas y contempla fases de selección de predicciones y simulación operativa con distintas modalidades
(long/short, scale-in, sizing dinámico). Las métricas evaluadas incluyen Sharpe, Max Drawdown, Calmar y CAGR; además se usan simulaciones Monte Carlo
para analizar robustez y sensibilidad. Para despliegue se dockeriza el sistema e integra monitorización con Prometheus/Loki, visualización en Grafana
y alertas vía Slack. El trabajo incluye análisis ético y económico, resultados de entrenamiento y backtests, y propone líneas futuras como aprendizaje
no supervisado, orquestación MLOps y optimización con Pyomo.

\section*{Palabras clave}

Trading algorítmico, Swing Charts, Machine Learning, LSTM, GRU, Triple barrera, Walk-forward, Feature engineering, Backtesting, Monte Carlo, RobustScaler, Dockerización, Grafana, Gestión de riesgo, Selección de características (LASSO)